% Figure | used in parts/03-integration/03-stokes/03-proof.tex
%!TEX root = ../main.tex
\begin{figure}[ht]
\centering
\begin{tikzpicture}[>=Stealth,scale=1.05]
  \fill[blue!12] (0,0) circle (1.9);
  \fill[white] (0,0) circle (0.65);
  \draw[navy] (0,0) circle (1.9);
  \draw[navy] (0,0) circle (0.65);
  \node[navy] at (-1.15,0.55) {$M$};
  % outer boundary: ccw arrow
  \draw[->,navy,thick] (60:1.9) arc (60:75:1.9);
  \draw[->,red,thick] (45:1.9)--(45:2.75) node[above right] {$N$};
  % inner boundary: clockwise arrow, N points into the hole
  \draw[->,navy,thick] (200:0.65) arc (200:185:0.65);
  \draw[->,red,thick] (330:0.65)--(330:0.15);
\end{tikzpicture}

\small $N=$ outward normal of $M$ (on the inner circle it points into the hole)
\caption{An annulus $M$ with its induced boundary orientation: the outer circle is traversed counter-clockwise and the inner circle clockwise, so that the region lies to the left. Applied to the angle form of Example~\ref{ex:angle} ($\dd\omega=0$) Stokes gives $\oint_{\mathrm{outer}}\omega-\oint_{\mathrm{inner\ (ccw)}}\omega=0$: both equal $2\pi$.}
\label{fig:annulus}
\end{figure}
